% Options for packages loaded elsewhere
\PassOptionsToPackage{unicode}{hyperref}
\PassOptionsToPackage{hyphens}{url}
\documentclass[14pt]{extarticle}
\usepackage{extsizes}
\usepackage{xcolor}
\usepackage{amsmath,amssymb}
\setcounter{secnumdepth}{-\maxdimen} % remove section numbering
\usepackage{iftex}
\ifPDFTeX
  \usepackage[T1]{fontenc}
  \usepackage[utf8]{inputenc}
  \usepackage{textcomp} % provide euro and other symbols
\else % if luatex or xetex
  \usepackage{unicode-math} % this also loads fontspec
  \defaultfontfeatures{Scale=MatchLowercase}
  \defaultfontfeatures[\rmfamily]{Ligatures=TeX,Scale=1}
\fi
\usepackage{lmodern}
\ifPDFTeX\else
  % xetex/luatex font selection
\fi
% Use upquote if available, for straight quotes in verbatim environments
\IfFileExists{upquote.sty}{\usepackage{upquote}}{}
\IfFileExists{microtype.sty}{% use microtype if available
  \usepackage[]{microtype}
  \UseMicrotypeSet[protrusion]{basicmath} % disable protrusion for tt fonts
}{}
\makeatletter
\@ifundefined{KOMAClassName}{% if non-KOMA class
  \IfFileExists{parskip.sty}{%
    \usepackage{parskip}
  }{% else
    \setlength{\parindent}{0pt}
    \setlength{\parskip}{6pt plus 2pt minus 1pt}}
}{% if KOMA class
  \KOMAoptions{parskip=half}}
\makeatother
\setlength{\emergencystretch}{3em} % prevent overfull lines
\providecommand{\tightlist}{%
  \setlength{\itemsep}{0pt}\setlength{\parskip}{0pt}}
\usepackage{bookmark}
\IfFileExists{xurl.sty}{\usepackage{xurl}}{} % add URL line breaks if available
\urlstyle{same}
\hypersetup{
  hidelinks,
  pdfcreator={LaTeX via pandoc}}

\author{}
\date{}

\begin{document}

\textbf{Base de Datos}

\textbf{Daniel Rodulfo, Víctor Pinto}

\begin{itemize}
\item
  \textbf{Introducción}
\end{itemize}

\begin{quote}
En el presente documento, abordaremos la definición informática de lo
que es una base de datos. Comenzaremos con los conceptos mas básicos, su
evolución y el estado actual de esta rama de la computación, hasta
establecer una lista de componentes con el cual armar una base de datos
de propósito general con un presupuesto delimitado.
\end{quote}

\begin{itemize}
\item
  \textbf{Que es una base de datos}
\end{itemize}

\begin{quote}
Una base de datos se define como un conjunto de datos estructurados de
tal manera que puedan ser accedidos con facilidad y rapidez, asi como
gestionados o consultados por un equipo informático.

Una base de datos debe contar con una organización que no comprometa su
seguridad, permita la búsqueda eficiente y la recuperación de la
información.

La base de datos por si sola se encarga unicamente de almacenar la
información, con lo cual se requiere de un programa, ajeno a la base en
si, cuya tarea consista en gestionar como se almacenan los datos,
quienes pueden acceder a ella o a datos específicos, permisos y las
operaciones que se pueden aplicar sobre los datos.
\textbackslash\textbackslash{}

El software encargado de administrar y regular la base de datos se le
conoce como database management system (DBMS). Este posee un repertorio
de opciones básicas con las cuales puede alterar una base de datos en
concreto: crear, leer, actualizar y eliminar datos. La nomenclatura de
estas operaciones dependera del lenguaje de consulta que se emplee.

Las bases de datos son ampliamente utilizadas en toda clase de sectores
en donde se necesite almacenar información y disponer de esta con
celeridad. Desde sistemas bancarios, analisis de datos y nube,
aplicaciones moviles e industriales, entre otros.

Debido a que la manera en que una base de datos almacena su información
puede variar segun lo deseado, existen diferentes tipos de bases de
datos, como las jerárquicas, las de objetos, relacionales, y una
variedad de categorias que se ajustan al como accedemos a la información
de la base de datos y el propósito de
esta.\textbackslash\textbackslash{}

Dentro de una base de datos, la estructuración no es lo único que puede
cambiar, sino tambien los datos que la base almacena. Si estos datos
disponen de un formato especifico y definido, estaremos hablando
entonces de datos estructurados, los cuales se almacenan dentro de una
base de datos relacional. Estos formatos, que pueden ser fechas,
números, caracteres y cadenas, posibilita la relación entre los datos,
esto es, que existe una similitud en formato de los datos dentro de la
base de datos. Cuando los datos carecen de un formato común y
especifico, estaremos tratando con datos datos no estructurados, que se
guardaran en bases de datos no relacionales. Estos pueden ser
documentos, ficheros y multimedia, y suelen tratarse con NoSQL.

\textbf{Bases de Datos Relacionales}: En una base de datos relacional,
los datos se almacenan en filas y columnas en tablas diferenciadas,
pudiendo acceder a los datos mediante algoritmos de búsqueda que
identifican etiquetas de relación únicas, las llamadas claves.

\textbf{Lenguaje de consulta SQL:} Las bases de datos relacionales
utilizan un modelo de consulta estructurado, tambien llamado SQL, para
buscar, alterar o eliminar la información. Un lenguaje de base de datos
puede ser de definición o de manipulación de los datos. Los lenguajes de
manipulación incluyen los lenguajes procedurales, que permiten utilizar
iteraciones o recursividad. SQL es un lenguaje de definición, pero puede
integrarse dentro de lenguajes host procedurales, como C++, Java o
Python.

\textbf{Bases de datos no relacionales:} Se utilizan para manejar datos
masivos de manera dinámica, sin usar el formato de fila y columna.
Cuentan con una gran escalabilidad y pueden emplear esquemas mas
flexibles.

\textbf{Características} de las Bases de datos:

Las características de una base de datos puede diferir de otra, aun asi
existen pautas que siempre van a ser convenientes para el diseño de una
base de datos.

\textbf{Consistencia:} La consistencia se basa en el modelo ACID y el
modelo BASE, los cuales aseguran que los datos se guardaran de forma
coherente y segura, sin que se corrompan en ningún modo

\textbf{Normalización:} Consiste en que, al momento de organizar la
información en una base de datos, se crean tambien tablas, relaciones y
reglas para proteger los datos y evitar la redundancia de estos.

\textbf{Persistencia:} La persistencia garantiza que cualquier
información pueda recuperarse incluso aunque el sistema que la estaba
consultado se desconecta.

\textbf{Escalabilidad:} Una base de datos debe mantenerse estable aun
cuando sus datos crezcan enormemente, siendo capaz de gestionar un mayor
volumen de datos de manera coherente

\textbf{Integridad:} La seguridad de los datos, su coherencia y su
precisión deben garantizarse para formar una buena base de datos

\textbf{Cardinalidad:} Se refiere a la relación que tiene una tabla
dentro de una base de datos con otra tabla. Por ejemplo, si es de muchos
a muchos, de uno a uno, o cual relación pueda describirse con otros
datos
\end{quote}

\begin{itemize}
\item
  \textbf{Breve historia de las bases de datos}
\end{itemize}

\begin{quote}
Durante los 60, la gestión de los datos de manera electrónica se fue
volviendo cada vez mas compleja, pues se debía supervisar y repartir
permisos constantemente. Fue entonces cuando se comenzó a divisar una el
concepto de una capa separada del software y el sistema operativo.
Comenzaron a implementarse modelos de bases de datos jerárquicos y de
red, que eran poco flexibles y sumamente simples. IBM revoluciono la
industria al introducir un sistema de base de datos relacional en los
70s.

Posteriormente, durante los 80s, llegaron al mercado los lenguajes de
consulta de bases de datos como SQL y los sucesores de IBM.

Para los 2000s se popularizaron los lenguajes NoSQL junto a proyectos de
codigo libre que aportaban sus propias contribuciones al mercado, como
MySQL y PostgreSQL.

Evolución de los Modelos de bases de datos:

\textbf{Jerárquico}: Se trata del modelo mas antiguo. En este modelo,
hay una escala del tipo árbol, en donde los nodos padres engendran nodos
hijos, ordenados en jerarquía. Los niveles de una jerarquía no estan
relacionados, y los nodos hijos siempre debian derivar de un nodo padre.
Es un modelo inflexible por su estructura, a cambio de ser facil de
identificar

\textbf{Red}: Desarrollados a la par de los modelos relacionales, en
este modelo los nodos no tienen una relación definida ni limitada con
sus pares. Cada registro puede tener diferentes precedentes, por lo que
no existe un único camino para llegar a un registro en particular.
Pueden eliminarse nodos sin que esto afecte directamente a la totalidad
de la red

\textbf{Relacional:} Se trata del modelo mas generalizado y popular en
la actualidad. Usa como lenguaje de consulta el SQL, y su modelado se
basa en el empleo de tablas y relaciones. Las tablas determinan la
ubicación de cada dato; cada dato conforma un registro (fila) y se
guardan como atributos (columna). Se utilizan claves para acceder a los
datos

\textbf{Orientado a objetos:} Aparecen en los 80s pero con escasa
aplicabilidad. Se guardan en los objetos los métodos necesarios para
acceder a la información

\textbf{Orientado a documentos}: Los documentos son la unidad de básica
de almacenamiento. Cada dato, guardados en documentos individuales,
consta de una clave y de un valor. No se define una cantidad de pares
clave valor, lo que genera que la disparidad entre documentos exista
\end{quote}

\begin{itemize}
\item
  \textbf{Componentes básicos}
\end{itemize}

\begin{quote}
Los componentes mas básicos requeridos para soportar una base de datos
son los siguientes:

CPU: El cerebro de la maquina. Se encarga de gestionar las consultas y
la transacción de los datos

Memoria RAM: Fundamental para la rapidez de los datos, al mantener en
cache las consultas mas frecuentes o las relaciones mas predominantes.

Almacenamiento: Unidades físicas donde se almacenan los datos para su
persistencia. Para las tareas de consulta de una base de datos, se
prefiere unidades del tipo SSD para las lecturas rapidas

Hardware de Red: Equipo que permita comunicarse con la red para las
comunicaciones al servidor de la base de datos.
\end{quote}

\begin{itemize}
\item
  \textbf{Estado actual del arte}
\end{itemize}

\begin{quote}
La computación se encuentra en constante cambio y evolución. Las
novedades que conciernen a la base de datos en la actualidad estan
profundamente relacionadas con la llegada y el auge de la inteligencia
artificial. La tendencia marca que las bases de datos estan dejando de
ser simples repositorios de información, para pasar a ser plataformas
inteligentes y autónomas, capaces de hacer transacciones con la IA.

Funcionalidades como la búsqueda vectorial comienzan a ser nativas en
diferentes arquitecturas de bases de datos, sin requerir
especialización. Tambien comienzan a difuminarse las fronteras entre
bases de datos exclusivas para transacciones y analisis.

Aunado a esto, las bases de datos pueden moldearse para ser aptas para
agentes de IA que escriben y ejecutan codigo, con un catalogo semántico
y un espacio seguro para el codigo creado.

Se ve una marcada consolidación del PostgreSQL como base de datos de
referencia para las IAs, debido a su codigo abierto y flexibilidad
\end{quote}

\begin{itemize}
\item
  \textbf{Tabla con los componentes y presupuesto}
\end{itemize}

1-Placa madre: ASUS ProArt X870E-CREATOR WiFi \$494,99.

Motivos:

\begin{itemize}
\item
  PCIe 5.0 en 2 slots M.2: Mayor velocidad de lectura/escritura, menor
  latencia y más IOPS en el almacenamiento NVMe.
\item
  soporte ECC: Para futuras actualizaciones.
\item
  VRM de 16+2+2 fases: Garantizas entrega de energía estable al CPU bajo
  carga sostenida 24/7.
\item
  4 slots ram- permite llegar 256 gb de DDR5.
\end{itemize}

link:https://www.amazon.com/-/es/ProArt-X870E-CREATOR-completo-generaci\%C3\%B3n-potencia/dp/B0DF123GCV

2- Cpu: AMD Ryzen 9 9950X3d \$639.99.

Motivos:

- Memoria l3 de 144 mb: Permite mantener mas índices calientes con el
fin de incrementar la velocidad de respuesta.

- 16 núcleos / 32 hilos: Alta concurrencia, múltiples consultas
simultáneas sin colas.

-Soporte ECC: preparado para RAM ECC en una futura actualización.

link:https://www.amazon.com/-/es/AMD-Ryzen-9950X3D-procesador-n\%C3\%BAcleos/dp/B0DVZSG8D5

3- ram: G.Skill Trident Z5 RGB 64GB DDR5 6400 CL32 \$858.

Motivos:\\
64GB DDR5 --- mayor capacidad en RAM significa menos accesos al disco,
queries más rápidas.

6400MHz --- mayor frecuencia que otros kits considerados, más ancho de
banda para el motor de BD.

CL32 --- baja latencia, respuesta más rápida en accesos aleatorios
frecuentes.

AMD EXPO --- perfil nativo para AM5, rendimiento máximo garantizado sin
configuración manual.

Disipador incluido --- temperatura estable bajo carga sostenida 24/7.

link:https://www.amazon.com/-/es/G-Skill-Trident-6400MHz-1-40V-doble/dp/B0BSP5F4CN

4-NVMe OS: Samsung 990 Pro 1TB \$249,99.

Motivos:

\begin{itemize}
\item
  Disco dedicado exclusivo al OS: Evita competencia por el canal de I/O
  con la BD.
\item
  Aislamiento de fallos: Si un disco falla el otro no se ve afectado.
\item
  PCIe Gen4 suficiente: el OS no se beneficia de la velocidad de gen5.
\end{itemize}

link:https://www.amazon.com/-/es/Samsung-velocidades-secuencial-computaci\%C3\%B3n-estaciones/dp/B0BHJF2VRN

5-NVMe BS:WD Black SN8100 2TB Gen5 \$430,28

Motivos:

\begin{itemize}
\item
  PCIe Gen5: Máxima velocidad de lectura/escritura para acceso constante
  a datos.
\item
  2.3 millones de IOPS: Operaciones aleatorias de alta velocidad.
\item
  2TB de capacidad.
\item
  Heatsink: Temperatura estable bajo carga sostenida.
\end{itemize}

link:https://www.amazon.com/-/es/WD\_Black-disipador-velocidad-escritura-aplicaciones/dp/B0F3H14JQG

6-Fuente de poder: be quiet! Power Zone 2 750W 80+ Platinum \$99.90

\begin{itemize}
\item
  80+ Platinum: Mínimo desperdicio de energía en la conversión de
  corriente alterna (AC) a corriente continua (DC),menos calor y menor
  consumo eléctrico para carga sostenida.
\item
  750W:margen suficiente para alimentar a todos los componentes.
\item
  Zero RPM: Ventilador silencioso a baja carga, ideal para operación
  continua.
\item
  ATX 3.1 / PCIe 5.1: compatibilidad con estándares modernos para
  futuras actualizaciones.
\end{itemize}

link:
https://www.amazon.com/-/es/quiet-alimentaci\%C3\%B3n-Eficiencia-Cybenetics-BP006US/dp/B0DVJGSK2W

7-Refrigeración:Noctua NH-D15S \$104.95

Motivos:

\begin{itemize}
\item
  Refrigeración por aire: Mas confiable a largo plazo que la
  refrigeración liquida para carga sostenida.
\item
  Doble torre: Mayor superficie de disipación para manejar los 170W TDP
  del 9950X3D bajo carga sostenida.
\item
  Diseño: 65mm de espacio libre para la RAM, evita conflictos de
  compatibilidad física con los módulos DDR5.
\end{itemize}

link:https://www.amazon.com/-/es/Noctua-NH-D15S-Enfriador-premium-ventilador/dp/B00XUVGLEU

8-Carcasa: Lian Li LANCOOL 217 \$119,99.

Motivos:

\begin{itemize}
\item
  5 ventiladores pre-instalados --- 2x170mm frontales, 2x120mm
  laterales, 1x140mm trasero, para una correcta cobertura de aire.
\item
  Diseño orientado al flujo de aire: Ideal para operación sostenida.
\item
  Espacio para expansión: instalación flexible de PSU y cables dejando
  espacio para discos adicionales.
\end{itemize}

link:https://www.amazon.com/dp/B0DWF95QP7

Totalizacion:\$494.99 + \$639.99 + \$858.00 + \$249.99 + \$430.28 + \$99.90 + \$104.95 + \$119.99 = \$2,998.09

\end{document}
